%% kapa: exam=1 round=21 year=2010 period=2 form=A subject=appraisal_law subject_ko=부동산관계법규 seq=39 number=79
%% source: Q-Net 감정평가사 기출문제 (cst00309) artlSeq=5024181 file=extracted/1008/5024181/2010년 제21회 감정평가사 1차시험 기출문제/2교시.hwp
%% conversion: kordoc 4.24.51 -> tools/gen1.py (auto); review=proofread-gpt-sol; flags=none
\begin{question}{79}
부동산등기법령상 등기관이 지적공부 소관청 또는 건축물대장 소관청에 지체 없이 그 사실을 알려야 할 경우가 아닌 것은?
\begin{choices}
  \item 소유권의 말소등기를 한 경우
  \item 소유권의 등기명의인 표시의 변경등기를 한 경우
  \item 동일한 토지에 중복하여 등기된 것을 발견한 경우
  \item 소유권의 보존등기를 한 경우
  \item 소유권의 경정등기를 한 경우
\end{choices}
\end{question}
\answer{3}
